\chapter{Termoidraulica del reattore \label{cap3}}


\section{Equazioni di conservazione}
L'equazione di conservazione della massa, della quantità di moto e 
dell'energia 
possono essere ricavate a partire dalla formulazione dell'equazione integrale 
di 
conservazione per una grandezza intensiva.
La descrizione classica del moto di un fluido si basa sull'ipotesi di mezzo 
continuo, in cui è possibile considerare il moto di elementi infinitesimali di 
fluido.
La velocità istantanea dell'elemento infinitesimale di fluido che occupa la 
posizione $(x,y,x)$ al tempo $t$ è definita dal campo di velocità 
$\vec{\upsilon}(x,y,x,t)$.
Inoltre se si assume l'elemento come uno sistema termodinamico chiuso, è 
possibile associare uno stato termodinamico definito dalle variabili di stato 
di 
temperatura $T$, pressione $p$, e di densità $\rho$.


Si consideri una grandezza estensiva $\Psi$ compresa in una regione di spazio 
tridimensionale $\Omega$, delimitata da una superficie chiusa $\partial \Omega$.
In termodinamica una grandezza estensiva è tale in quanto dipende dalla 
quantità 
di materia.
Quindi si può definire $\psi$, una grandezza intensiva, che corrisponde alla 
grandezza estensiva $\Psi$ per unità di massa.
Inoltre si definisce la densità di flusso $\vec{J}$ che descrive l'interazione 
di $\psi$ con il sistema, la densità del fluido $\rho$, e il termine sorgente 
$\mathcal{S}$ per unità di tempo e di massa.
Si può formulare nella forma più generale l'equazione integrale di 
conservazione 
della grandezza intensiva $\psi$ sul volume $\Omega$, che esprime l'evoluzione 
temporale di tale grandezza,
\begin{equation}
\label{eq_integrale_di_conservazione}
\frac{d}{dt}\int_{\Omega}\rho\psi \, dV = - \int_{\partial \Omega} \vec{J} 
\cdot 
\hat{n}\, dS + \int_{\Omega} \rho \mathcal{S} \, dV .
\end{equation}
Dall'equazione \eqref{eq_integrale_di_conservazione} si desume che la variazione 
della grandezza intensiva su tutto il volume è uguale al bilancio tra il flusso 
che attraversa la superficie $\partial \Omega$ e il contributo della sorgente.
Nel caso le funzione sia sufficientemente regolari è possibile applicare il 
teorema del trasporto di Reynolds, ottenendo
\begin{equation}
\label{eq_integrale_reynolds}
\frac{d}{dt}\int_{\Omega}\rho\psi \, dV = - \int_{\partial \Omega} 
(\rho\psi\vec{\upsilon}) \cdot \hat{n}\, dS + \int_{\Omega} \frac{\partial 
\rho\psi }{\partial t} \, dV ,
\end{equation}
dove $\vec{\upsilon}$ rappresenta il vettore velocità istantanea dell'elemento 
infinitesimale che attraversa la superficie $\partial \Omega$.
Il teorema di Gauss-Green permette di trasformare gli integrali di superficie 
in 
integrali di volume, dove si ottiene rispettivamente
\begin{equation}
\int_{\partial \Omega} \vec{J} \cdot \hat{n}\, dS = \int_{\Omega} \nabla \cdot 
\vec{J} \, dV \,,
\end{equation}

\begin{equation}
\label{greengauss}
\int_{\partial \Omega} (\rho\psi\vec{\upsilon}) \cdot \hat{n}\, dS = 
\int_{\Omega} \nabla \cdot (\rho\psi\vec{\upsilon})  \, dV \,.
\end{equation}

Se si introducono le equazioni \eqref{eq_integrale_reynolds}-\eqref{greengauss} 
nella equazione \eqref{eq_integrale_di_conservazione}, si ottiene
\begin{equation}
\label{eq_termo}
 \int_{\Omega} \frac{\partial \rho\psi }{\partial t} \, dV + \int_{\Omega} 
\nabla \cdot (\rho\psi\vec{\upsilon})  \, dV = -\int_{\Omega} \nabla \cdot 
\vec{J} \, dV + \int_{\Omega}\rho \mathcal{S} \,dV \,.
\end{equation}
L'equazione \eqref{eq_termo} deve essere verificata su tutti i volumi $\Omega$, 
quindi è possibile rimuovere l'integrale, ottenendo l'equazione differenziale
\begin{equation}
\label{eq_converse_diff}
 \frac{\partial \rho\psi }{\partial t}  +  \nabla \cdot 
(\rho\psi\vec{\upsilon}) 
  = -\nabla \cdot \vec{J}+ \rho \mathcal{S} \,.
\end{equation}
L'equazione di conservazione in forma differenziale può essere scritta per le 
diverse grandezze ottenendo così le equazione di conservazione della massa, 
della quantità di moto e dell'energia.

\subsection{Equazione di conservazione della massa}
Il principio di conservazione della massa stabilisce che la massa non può 
essere 
creata o distrutta. L'equazione di conservazione della massa si ottiene 
sostituendo all'equazione \eqref{eq_converse_diff} i termini
\begin{equation}
\psi=1 \qquad \vec{J}=0 \qquad \mathcal{S}=0 \,.
\end{equation}
Si ottiene così l'equazione di continuità alle derivati parziali,
\begin{equation}
\label{eq_continu}
\frac{\partial \rho}{\partial t} + \nabla \cdot (\rho \vec{\upsilon}) = 0 \,.
\end{equation}
Inoltre se si assume l'ipotesi di fluido incomprimibile, ovvero considerando la 
densità costante $\rho(x,y,z,t)=cost$, l'equazione \eqref{eq_continu} si 
semplifica 
\begin{equation}
\label{..}
\nabla \cdot \vec{\upsilon} = 0 \,,
\end{equation}
che esprime come attraverso una qualunque superficie la quantità di massa 
entrante è pari a quella uscente, in quanto si è considerato il fluido 
incomprimibile.

\subsection{Equazione di conservazione della quantità di moto}
L'equazione di bilancio della quantità di moto si ottiene uguagliando la 
variazione delle quantità di moto del sistema alle forze applicate sul volume 
d'integrazione $\Omega$.
Le forze che agiscono sul sistema si dividono in in forze volumetriche e in 
forze di superficie.
Le forze volumetriche agiscono sul volume del fluido e sono dovute a campi 
esterni, come quello gravitazionale, elettrico o magnetico.
\'E stato assunto che il fluido considerato sia elettricamente neutro ed 
immerso 
in un campo gravitazionale così che l'unica forza volumetrica sia il forza di 
gravità $\vec{g}$.
Le forze di superficie agiscono sulla superficie dell'elemento infinitesimale 
di 
fluido e sono espresse dal tensore delle tensioni $\vec{\sigma}$, suddiviso in 
tensore idrostatico e in tensore delle deformazioni.
Sostituendo all'equazione \eqref{eq_integrale_di_conservazione} le seguenti 
grandezze
\begin{equation}
\psi= \vec{\upsilon} \qquad \vec{J}=-\vec{\sigma}= p \vec{I} - \vec{\tau} 
\qquad 
\mathcal{S}=\vec{g} \,,
\end{equation}
dove il tensore delle tensioni $\vec{\sigma}$ è stato scomposto nel termine di 
pressione $p\vec{I}$, dove $\vec{I}$ è il tensore identità, e nel tensore degli 
sforzi viscosi $\vec{\tau}$. si ottiene l'equazione di Navier-Stokes.
Infatti,
\begin{equation}
\label{eq_navier-stokes}
\frac{\partial \rho \vec{\upsilon}}{\partial t} + \nabla \cdot (\rho 
\vec{\upsilon}\vec{\upsilon})= -\nabla P + \nabla \cdot \vec{\tau} + \rho 
\vec{g} \,.
\end{equation}

Per un fluido newtoniano, ovvero un fluido la cui viscosità rimane costante al 
variare della velocità, il tensore degli sforzi viscosi è espresso dalla 
relazione,
\begin{equation}
\vec{\tau}=\mu [(\nabla\vec{\upsilon})+(\nabla\vec{\upsilon})^T]  
\biggl(\frac{2}{3}\mu-\lambda\biggr)(\nabla \cdot \vec{\upsilon})\vec{I} \,,
\end{equation}
dove $\mu$ è la viscosità dinamica e $\lambda$ la viscosità di massa.

Assumendo l'ipotesi di fluido incomprimibile, ovvero che dove la densità del 
fluido si può considerare costante, si può semplificare il secondo termine che 
compare nel tensore degli sforzi. Questo può essere fatto in quanto vale l'equazione di continuità 
\eqref{eq_continu}.
Inoltre se si considera la viscosità costante, vale la seguente relazione,
\begin{equation}
\nabla \cdot \vec{\tau} = \mu \nabla^2 \vec{\upsilon}
\end{equation}
Attraverso queste ipotesi appena fatte e semplificando il termine di avvezione 
$\nabla \cdot (\rho 
\vec{\upsilon}\vec{\upsilon})=\vec{\upsilon}\cdot\nabla\vec{\upsilon}$, è 
possibile semplificare l'equazione \eqref{eq_navier-stokes}
\begin{equation}
\frac{\partial \vec{\upsilon}}{\partial t} + 
\vec{\upsilon}\cdot\nabla\vec{\upsilon} = -\frac{1}{\rho} \nabla p + \mu 
\nabla^2\vec{\upsilon}+\vec{g} \,.
\end{equation}

\subsection{Equazione di conservazione dell'energia}
L'equazione di conservazione dell'energia rappresenta la traduzione matematica 
del primo principio della termodinamica. 
La grandezza intensiva è la somma dell'energia interna e dell'energia cinetica 
dell'elemento infinitesimale di fluido. Infatti sostituendo all'equazione 
\eqref{eq_integrale_di_conservazione} le seguenti grandezze
\begin{equation}
\psi= \varepsilon+\frac{|\vec{\upsilon}|^2}{2} \qquad 
\vec{J}=\vec{q}-\vec{\sigma}\cdot \vec{\upsilon} \qquad 
\mathcal{S}=\frac{\dot{Q}}{\rho}+\vec{g}\cdot\vec{\upsilon} \,,
\end{equation}
dove $\varepsilon$ rappresenta l'energia interna, $\frac{u^2}{2}$ l'energia 
cinetica per unità di massa, $\vec{q}$ il flusso di calore che attraversa la 
superficie $\partial \Omega$ e $\dot{Q}$ la sorgente volumetrica di calore.
Si ottiene l'equazione di conservazione dell'energia sostituendo all'equazione 
\eqref{eq_integrale_di_conservazione} e moltiplicando scalarmente l'equazione 
con l'equazione di conservazione della quantità di moto e sottraendo membro a 
membro si ottiene
\begin{equation}
\label{eq_conse_energia}
 \frac{\partial}{\partial t}(\rho \varepsilon) + \nabla \cdot (\rho \varepsilon 
\vec{\upsilon}) = -\nabla \cdot \vec{q} - p \nabla \cdot \vec{\upsilon} + 
\vec{\tau} : \nabla \vec{\upsilon}+ \rho \vec{g}\cdot \vec{\upsilon} + \dot{Q} 
\, .
\end{equation}
Aggiungendo l'ipotesi di fluido incomprimibile, quindi considerando valida 
l'equazione di continuità \eqref{eq_continu}, e utilizzando la definizione di 
derivata sostanziale, si ottiene una forma più compatta dell'equazione 
\eqref{eq_conse_energia} nella forma
\begin{equation}
 \rho \frac{d\varepsilon}{dt}= - \nabla \cdot \vec{q} + 
\vec{\tau}:\nabla\vec{\upsilon}+ \rho \vec{g}\cdot \vec{\upsilon} + \dot{Q}\,.
\end{equation}
Si introduce la legge di Fourier, che correla il flusso di calore al campo di 
temperatura
\begin{equation}
 \vec{q} = - \kappa(T)\nabla T \,,
\end{equation}
dove $\kappa$ rappresenta la conducibilità termica del fluido.
Inoltre si pone 
\begin{equation}
 \vec{\tau} : \nabla \vec{\upsilon} = \mu\Phi \,,
\end{equation}
dove $\Phi$ è la funzione di dissipazione viscosa e $\mu$ il coefficiente di 
viscosità.
Nell'ipotesi di fluido incomprimibile, l'energia interna $\varepsilon$ è in 
funzione della sola temperatura
\begin{equation}
 d\varepsilon=c_v(T)dT
\end{equation}
dove $c_v(T)$ rappresenta il calore specifico a volume costante del fluido 
considerato.
A seguito delle approssimazioni precedenti si ottiene
\begin{equation}\label{eq dell'energia}
 \rho c_v\frac{dT}{dt}=\nabla \cdot (\kappa \nabla T) + \mu \Phi + \rho 
\vec{g}\cdot \vec{\upsilon} + \dot{Q} \,.
\end{equation}

\section{Modello del reattore}

\subsection{Modello del core e zone ingresso}
In un reattore nucleare sono presenti una zona di ingresso 
e una di uscita dette comunemente plena, ovvero zone in cui si raccoglie 
il refrigerante: il plenum superiore e il plenum inferiore localizzati 
rispettivamente sopra e sotto il nocciolo. Il refrigerante attraversa dal basso 
verso l'altro, le assembly di combustibile, in cui esso asportata 
il calore 
generato dalle reazioni di fissione.
Infatti in queste regioni, il refrigerante fluisce in un dominio 
tridimensionale 
aperto e lo stato del refrigerante, definito dal campo di velocità 
$\vec{\upsilon}$, pressione $p$ e temperatura $T$, può essere determinato dalle 
equazioni seguenti di conservazione tridimensionali per un fluido 
incompressibile 
per un modello 3D-CFD
\begin{equation}
 \nabla \cdot (\rho \vec{\upsilon}) = 0 \,,
\end{equation}
\begin{equation}
 \frac{\partial \rho \vec{\upsilon}}{\partial t} + \nabla \cdot 
(\rho\vec{\upsilon}\vec{\upsilon}) = - \nabla p + \nabla : \vec{\tau} \rho 
\vec{g} \,,
\end{equation}
\begin{equation}
\frac{\partial \rho h}{\partial t} + \nabla \cdot (\rho \vec{\upsilon} h) = 
\Phi 
+ \nabla \cdot (\frac{\kappa_f}{c_p} \nabla h) + \dot{Q} \,,
\end{equation}
dove $c_p$ è il calore specifico a pressione costante, $\kappa_f$ la 
conduttività termica effettiva e $h$ l'entalpia.
Quest'ultima è in funzione della temperatura ed è definita dall'equazione,
\begin{equation}
 h(T)=c_p T \,,
\end{equation}
considerando il $c_p$ costante all'interno della regione considerata.
In particolare i termini che più caratterizzano le proprietà del refrigerante, 
come il piombo, in queste regioni sono la viscosità del fluido $\mu_f$, 
la conducibilità termica $\kappa_f$, la capacità termica specifica e la densità 
$\rho$, le quali possono essere considerate in funzione della sola
temperatura $T$.

In molte componenti del reattore la simulazione del flusso interno del 
refrigerante è molto complessa e pertanto è necessario introdurre 
approssimazioni, come l'omogeneizzazione del dominio.
Inoltre la regione del nocciolo è così complessa e dettagliata, a causa della 
presenta di diverse componenti sia solide che liquide, che un approssimazione 
ad 
un modello poroso diventa necessario in quanto una simulazione della 
distribuzione di velocità, pressione e temperatura non è possibile. 
\cite{cerroni}

La distribuzione delle perdite di pressione, definita dalle griglie presenti 
nel 
reattore e dalle assembly, può essere definita dal modulo di $|\beta|$ dal 
termine di caduta di pressione discreta
\begin{equation}
\label{eq_cad}
 \frac{1}{2} \beta(\vec{x}) \cdot \vec{\upsilon}|\vec{\upsilon}|\,.
\end{equation}
Se si assume il flusso essere in prima approssimazione verticale allora 
l'equazione \eqref{eq_cad} diventa
\begin{equation}
 \frac{1}{2} |\beta| \, |\vec{\upsilon}|^2 \,,
\end{equation}
Questo permette l'approssimazione solita al caso monodimensionale.

Tuttavia è possibile implementare equazioni e modelli di turbolenza per una più approfondita rappresentazione fisica inerente alla termoidraulica del reattore.
Per esempio nello studio in referenza \cite{leadmontecu} è presentata una 
simulazione numerica del flusso turbolento di un fluido LBE (lega eutettica 
piombo-bismuto) all'interno di un assembly di un reattore di IV$^{a}$ 
generazione, in cui è stato risolto il sistema di equazioni di 
Reynolds-Average-Navier-Stokes seguente
\begin{equation}
 \nabla \cdot \vec{\upsilon}=0 \,,
\end{equation}
\begin{equation}
 \rho \frac{\partial \vec{\upsilon}}{\partial t} + \rho (\vec{\upsilon} \cdot \nabla)\vec{\upsilon}= -\nabla p + \nabla \cdot \bigl[\mu (\nabla \vec{\upsilon} + \nabla \vec{\upsilon}^T)- \rho \bar{\vec{\upsilon}'\vec{\upsilon}} \bigr] + \rho \vec{g} \,,
\end{equation}
\begin{equation}
 \rho c_p \Bigl(\frac{\partial T}{\partial t} + (\vec{\upsilon} \cdot \nabla) T \Bigr) = \nabla \cdot \bigl[\lambda \nabla T - \rho c_p \bar{\vec{\upsilon}'T'}\bigr]+ Q \,.
\end{equation}
Il sistema di equazioni è chiuso, in questo caso, da un altro sistema di equazioni a quattro parametri logaritmici di turbolenza allo scopo di ottenere il tensore degli sforzi di Reynolds $\rho \bar{\vec{\upsilon}'\vec{\upsilon}}$ e il flusso di calore turbolento $\rho c_p \bar{\vec{\upsilon}'T'}$.

Infatti metalli liquidi sono classificati con un numero di Prandtl relativamente 
piccolo, in quanto sono caratterizzati da bassa viscosità $\nu$ e alti valori di 
diffusività termica $\alpha$, definito dall'equazione
\begin{equation}
 Pr = \frac{\nu}{\alpha} \simeq 0.01-0.025 \,,
\end{equation}
e pertanto è necessario introdurre modelli sulla turbolenza adattati a numeri di Prandtl così piccoli.


% \Figure[scale=0.8]{./figure/itex/tempeq.pdf}{ ssssss .}{tempeq}
% \Figure[scale=0.8]{./figure/itex/canale.pdf}{ sss .}{canale}

\subsection{Temperatura del combustibile}

\Figure[scale=0.8]{./figure/itex/dist_temp_comb.pdf}{Distribuzione del profilo 
di 
temperatura in un elemento di combustibile.}{profil_temp}

Dai valori mediati di temperatura ottenuti è possibile calcolare il profilo di 
temperatura all'interno di un barra di combustibile, con dimensioni $2a$ per 
l'elemento di combustibile e $b$ per la guaina di rivestimento.
Considerando una geometria piana e il problema stazionario, il campo di 
temperatura nell'elemento di combustibile è descrivibile attraverso 
un'equazione 
di Poisson monodimensionale alle derivate ordinarie
\begin{equation}
k_f \frac{d^2 T}{dx^2} + \dot{Q} = 0  \,,
\end{equation}
dove $\dot{Q}$ è la sorgente di calore volumetrica e $k_f$ è la conducibilità 
termica del combustibile.
Essendo un'equazione differenziale del secondo ordine alle derivate ordinarie e 
non omogenea è richiesta per la sua risoluzione le condizioni al contorno; la 
prima esprime che la temperatura è massima al centro della barra e le 
condizione 
di simmetria rispetto all'asse verticale.
Si ottiene la distribuzione nell'elemento di combustibile,
\begin{equation}
T(x)=T_f - \frac{\dot{q}}{2k_f}x^2 \,,
\end{equation}
dove $\dot{q}(x)=\dot{Q}x$ è il flusso di calore per unità di lunghezza.
Nella zona tra la guaina di rivestimento e il combustibile, detto comunemente 
{\it gap}
è presente un gas inerte, 
generalmente argon o elio.
Lo scambio termico è rappresentato attraverso un coefficiente di convezione 
equivalente, calcolato come $h_g=k_g/l_g$, dove $l_g$ è la dimensione del gap. 
Si ottiene così lo scambio termico tra la superficie dell'elemento di 
combustibile e la guaina di rivestimento come
\begin{equation}
 \dot{q}= h_g (T_a - T_2) \,,
\end{equation}
dove $T_a$ è la temperatura sulla superficie del combustibile e $T_2$ la 
temperatura sulla superficie interna della guaina.
Attraverso la legge di Fourier è possibile il flusso di calore che attraversa 
la 
superficie della guaina a contatto con il refrigerante e si ottiene
\begin{equation}
T_2 - T_c = \frac{b}{k_c} \dot{q}\,,
\end{equation}
dove $T_c$ è la temperatura della guaina sulla superficie a contatto con il 
refrigerante e $k_c$ la conducibilità termica della guaina.

Aggiungendo infine il flusso di calore al refrigerante attraverso un 
coefficiente di scambio equivalente
\begin{equation}
 \dot{q}= h_r (T_c - T_r) \,,
\end{equation}
si ottiene il profilo di temperatura della barra di combustibile
\begin{equation}
 T_f - T_r = \Delta T_f + \Delta T_g + \Delta T_c + \Delta T_r \,.
\end{equation}
Sostituendo i termini ottenuti in precedenza e raccogliendo il flusso di 
calore, 
si ottiene la temperatura del combustibile
\begin{equation}
 T_f = T_r + \dot{q} \bigl[ \frac{a}{2 k_f} + \frac{1}{h_g} + \frac{b}{k_c} + 
\frac{1}{h_r} \bigr] \,.
\end{equation}




 





